\documentclass[11pt]{article}
\usepackage[margin=1in]{geometry}
\usepackage{amsmath, amssymb, amsthm}
\usepackage{booktabs}
\usepackage{hyperref}
\usepackage{enumitem}

\newtheorem{proposition}{Proposition}
\newtheorem{remark}{Remark}

\title{Advantage Bias Under Selective Gradient Routing in GRPO}
\author{small-rl project notes}
\date{}

\begin{document}
\maketitle

\begin{abstract}
We analyze a systematic bias in the gradient signal received by retain and forget
adapters under gradient routing with GRPO. When reward-hacking (RH) samples are
filtered from the retain adapter's training signal, the remaining ``good'' samples
carry negative mean advantage, introducing a bias that pushes the retain adapter
away from desirable behavior. We derive this bias, show that baseline invariance
fails under selective sampling, and characterize severity as a function of the RH
detection rate. We then analyze the forget adapter's corresponding positive mean
bias and argue that it is beneficial but not uniquely principled. After examining
several corrective approaches --- advantage renormalization, conditioned penalties,
and manual advantage assignment --- we arrive at a simple fixed reward penalty as
the recommended design: it preserves GRPO's adaptive normalization properties
while giving the retain adapter an active avoidance signal.
\end{abstract}

%% ====================================================================
\section{Setup}
%% ====================================================================

\subsection{GRPO Advantage Computation}

Consider a prompt $s$ with $G$ generations $\{a_1, \ldots, a_G\}$ sampled from
the current policy $\pi_\theta$. GRPO computes per-sample advantages via
group normalization:
\begin{equation}\label{eq:grpo-advantage}
    \hat{A}_i = \frac{r_i - \bar{r}}{\sigma_r + \epsilon},
    \qquad
    \bar{r} = \frac{1}{G}\sum_{j=1}^{G} r_j,
    \qquad
    \sigma_r = \mathrm{std}(\{r_j\}_{j=1}^{G})
\end{equation}
By construction, $\sum_{i=1}^{G} \hat{A}_i = 0$ within each group.

The normalization in~\eqref{eq:grpo-advantage} provides three important
properties for stable training:
\begin{enumerate}[nosep]
    \item \textbf{Scale invariance}: gradient direction is independent of
        reward magnitude.
    \item \textbf{Adaptive temperature}: $1/\sigma_r$ acts as an inverse
        temperature --- tight reward distributions (model is confident) produce
        amplified advantages (stronger gradient), while spread distributions
        produce compressed advantages (gentler gradient).
    \item \textbf{Per-group calibration}: each prompt group is normalized
        independently, accommodating different reward landscapes across prompts.
\end{enumerate}

\subsection{Gradient Routing}

A reward-hacking detector partitions each group into \emph{bad} samples
$\mathcal{B} = \{i : \text{is\_rh}(a_i)\}$ with $|\mathcal{B}| = K$ and
\emph{good} samples $\mathcal{G} = \{1,\ldots,G\} \setminus \mathcal{B}$
with $|\mathcal{G}| = G - K$.

In the two-pass training step:
\begin{itemize}[nosep]
    \item \textbf{Pass 1} (retain adapter): gradients flow only through retain
        parameters, using samples in $\mathcal{G}$ (exclusive mode) or all
        samples with forget parameters hooked to zero.
    \item \textbf{Pass 2} (forget adapter): gradients flow only through forget
        parameters, using samples in $\mathcal{B}$, with retain parameters
        hooked to zero.
\end{itemize}

Crucially, both passes use advantages $\hat{A}_i$ computed from the full group
via~\eqref{eq:grpo-advantage}. The current implementation does not recompute
advantages for either sub-batch.

%% ====================================================================
\section{Retain Adapter: Negative Mean Advantage}\label{sec:retain}
%% ====================================================================

\subsection{The Bias}

Since RH-flagged samples have the highest rewards in the group (that is why they
were flagged), they carry positive advantages. The zero-sum property then implies:
\begin{equation}\label{eq:negative-mean}
    \bar{A}_{\mathcal{G}}
    = \frac{1}{G-K}\sum_{i \in \mathcal{G}} \hat{A}_i
    = \frac{-\sum_{i \in \mathcal{B}} \hat{A}_i}{G - K}
    < 0
\end{equation}

The retain adapter's gradient estimate is:
\begin{equation}
    \hat{g}_{\text{retain}}
    = \frac{1}{|\mathcal{G}|} \sum_{i \in \mathcal{G}}
        \hat{A}_i \, \nabla_\theta \log \pi_\theta(a_i \mid s)
\end{equation}

\subsection{Why Baseline Invariance Fails}

In standard REINFORCE, subtracting any constant baseline $b$ from the advantage
does not bias the gradient, because:
\begin{equation}
    \mathbb{E}_\pi\bigl[b \cdot \nabla_\theta \log \pi_\theta(a \mid s)\bigr]
    = b \cdot \nabla_\theta \sum_a \pi_\theta(a \mid s)
    = b \cdot \nabla_\theta 1
    = 0
\end{equation}

However, when we restrict the sum to good samples, we are implicitly multiplying
by the indicator $\mathbf{1}[\text{good}(a)]$, which depends on $a$:
\begin{equation}
    \mathbb{E}_\pi\bigl[\mathbf{1}[\text{good}(a)] \cdot b
        \cdot \nabla_\theta \log \pi_\theta(a \mid s)\bigr]
    = b \cdot \nabla_\theta \Pr_\pi[\text{good} \mid s]
    \neq 0
\end{equation}

The bias introduced by the full-group baseline under selective sampling is:
\begin{equation}\label{eq:bias}
    \boxed{%
    \text{bias} = \bar{A}_{\mathcal{G}}
        \cdot \nabla_\theta \Pr_\pi[\text{good} \mid s]}
\end{equation}

Since $\bar{A}_{\mathcal{G}} < 0$, this bias pushes the retain adapter toward
\emph{decreasing} the probability of generating good completions.

\subsection{Decomposition of the Gradient}

We can decompose the retain gradient into a biased mean component and an
unbiased relative-ranking component:
\begin{equation}\label{eq:decomposition}
    \hat{g}_{\text{retain}}
    = \underbrace{\bar{A}_{\mathcal{G}}
        \sum_{i \in \mathcal{G}} \nabla_\theta \log \pi(a_i \mid s)
    }_{\text{bias term } (\bar{A}_{\mathcal{G}} < 0)}
    + \underbrace{\sum_{i \in \mathcal{G}}
        \bigl(\hat{A}_i - \bar{A}_{\mathcal{G}}\bigr)
        \nabla_\theta \log \pi(a_i \mid s)
    }_{\text{relative ranking (zero mean)}}
\end{equation}

The relative-ranking term preserves the correct ordering among good samples
(push toward the best, away from the worst). The bias term is the problem.

\subsection{Severity Scaling}

The magnitude of $\bar{A}_{\mathcal{G}}$ depends on the RH detection rate
$K/G$. Since advantages are zero-mean over the full group with unit variance
(approximately), the bad samples' mean advantage is $\bar{A}_{\mathcal{B}}
\approx c > 0$ for some constant depending on the reward distribution. Then:
\begin{equation}
    \bar{A}_{\mathcal{G}} = \frac{-K \cdot \bar{A}_{\mathcal{B}}}{G - K}
\end{equation}

\begin{table}[h]
\centering
\begin{tabular}{@{}lcl@{}}
\toprule
$K/G$ & $\bar{A}_{\mathcal{G}}$ (approx.) & Severity \\
\midrule
$1/16$ & $-0.07$ & Negligible \\
$4/16$ (75th pct.\ detector) & $-0.33$ & Comparable to within-good variance \\
$8/16$ & $-1.0$ & Bias dominates signal \\
\bottomrule
\end{tabular}
\caption{Approximate bias magnitude for $G=16$ generations.}
\label{tab:severity}
\end{table}

\subsection{Dynamics Over Training}

The bias worsens over the course of training:
\begin{enumerate}[nosep]
    \item \textbf{Early}: The model rarely produces RH content.
        $K \approx 0$, bias is negligible, retain adapter learns normally.
    \item \textbf{Mid}: The forget adapter begins producing RH content.
        $K$ grows, bias grows, retain adapter starts being pushed away from
        good behavior.
    \item \textbf{Late}: Significant RH rate. Large $K$, strong negative bias.
        Retain adapter is actively degraded.
\end{enumerate}

This creates a pernicious feedback loop: the retain adapter's degradation
\emph{accelerates} as training succeeds at the forget task --- precisely when
clean adapter separation matters most.

\subsection{Consequence for Inference Ablation}

LoRA adapters are initialized at zero, so the initial policy equals the base
model. With negative mean advantage, the retain adapter learns corrections that
\emph{decrease} probability of good completions. At inference with
\texttt{retain\_only} (forget adapter ablated), the resulting model is
\textbf{worse than the base model} at the retain task.

Neither PPO clipping nor the KL penalty ($\beta \cdot \mathrm{KL}$) resolves
this. Clipping bounds the step size but preserves the gradient direction. The
KL penalty anchors toward the reference policy but does not correct the bias ---
it merely limits how far the retain adapter drifts in the wrong direction.

%% ====================================================================
\section{Forget Adapter: Positive Mean Advantage}\label{sec:forget}
%% ====================================================================

\subsection{The Bias Direction is Correct}

By an analogous decomposition, the forget adapter's gradient (in exclusive
routing, where it only sees bad samples) is:
\begin{equation}
    \hat{g}_{\text{forget}}
    = \underbrace{\bar{A}_{\mathcal{B}}
        \sum_{i \in \mathcal{B}} \nabla_\theta \log \pi(a_i \mid s)
    }_{\text{bias term } (\bar{A}_{\mathcal{B}} > 0)}
    + \underbrace{\sum_{i \in \mathcal{B}}
        \bigl(\hat{A}_i - \bar{A}_{\mathcal{B}}\bigr)
        \nabla_\theta \log \pi(a_i \mid s)
    }_{\text{relative ranking (zero mean)}}
\end{equation}

The positive bias pushes toward \emph{increasing} probability of bad
completions --- the desired behavior for the forget adapter.

\subsection{The Bias Magnitude is Not Principled}

While the bias direction is correct, the full-group baseline $\bar{r}$ is not the
variance-optimal baseline for the restricted (bad-samples-only) estimator. The
optimal baseline for the restricted objective would be:
\begin{equation}
    b^* = \frac{\mathbb{E}\bigl[\mathbf{1}[\text{bad}] \cdot r
        \cdot \|\nabla \log \pi\|^2\bigr]}
    {\mathbb{E}\bigl[\mathbf{1}[\text{bad}]
        \cdot \|\nabla \log \pi\|^2\bigr]}
\end{equation}
which depends on the action and is generally not equal to $\bar{r}$.

The full-group baseline also creates an \textbf{uncontrolled coupling} between
the two adapters' learning dynamics: as the retain task improves
($\bar{r}_{\mathcal{G}} \uparrow$), the group mean $\bar{r}$ rises, and the
forget adapter's advantages $\hat{A}_{\mathcal{B}}$ shrink. This coupling is
incidental (a side effect of the shared baseline), not designed.

Despite this, the full-group baseline is adequate in practice:
the bias direction is correct, the zero-variance edge case is handled well
(Section~\ref{sec:zero-var}), and the coupling is mild for typical $K/G$ ratios.
We therefore leave the forget adapter's advantages unchanged and focus corrective
effort on the retain adapter.

\subsection{The Zero-Variance Edge Case}\label{sec:zero-var}

When all bad samples achieve maximum reward (e.g., all trigger the RH detector
at ceiling), we have $r_i \approx r_j$ for all $i, j \in \mathcal{B}$, so:
\begin{itemize}[nosep]
    \item The relative-ranking term vanishes (no variance among bad samples).
    \item The gradient is entirely the bias term:
        $\hat{g}_{\text{forget}} \approx \bar{A}_{\mathcal{B}}
        \sum_{i \in \mathcal{B}} \nabla \log \pi(a_i \mid s)$.
    \item This is fine: ``increase probability of all these equally-rewarded
        bad completions.''
\end{itemize}

\begin{proposition}
Renormalizing forget-adapter advantages over only bad samples is harmful.
When $\sigma_{\mathcal{B}} \approx 0$:
\begin{equation}
    \hat{A}_i^{\mathrm{renorm}}
    = \frac{r_i - \bar{r}_{\mathcal{B}}}{\sigma_{\mathcal{B}} + \epsilon}
    \approx \frac{0}{\epsilon}
    \approx 0
    \quad \forall\, i \in \mathcal{B}
\end{equation}
The forget adapter receives zero gradient and learns nothing.
\end{proposition}

Even when bad-sample rewards do vary, renormalization removes the helpful
positive bias, replacing ``move toward all bad completions'' with
``prefer the most-bad over the least-bad.'' The former is more useful for
learning to capture RH behavior.

\subsection{Classic vs.\ Exclusive Routing}

In \textbf{classic} routing, the forget adapter receives gradients from
\emph{both} passes:
\begin{itemize}[nosep]
    \item Pass 1 (good samples, both adapters active): negative mean advantage
        $\to$ pushes forget adapter \emph{away from} good completions.
    \item Pass 2 (bad samples, retain hooked): positive mean advantage
        $\to$ pushes forget adapter \emph{toward} bad completions.
\end{itemize}

These signals are complementary: ``avoid good'' and ``pursue bad.''
In \textbf{exclusive} routing, the forget adapter only receives the
``pursue bad'' signal, missing the ``avoid good'' component. This suggests
classic mode produces a richer training signal for the forget adapter.

%% ====================================================================
\section{Corrective Approaches for the Retain Adapter}\label{sec:corrections}
%% ====================================================================

We consider four approaches in order of increasing sophistication, and argue
that the last --- a fixed reward penalty with standard GRPO normalization ---
is the best design.

\subsection{Approach 1: Advantage Renormalization}\label{sec:renorm}

Recompute advantages for retain-adapter training using only good samples:
\begin{equation}\label{eq:renorm}
    \hat{A}_i^{\text{retain}}
    = \frac{r_i - \bar{r}_{\mathcal{G}}}{\sigma_{\mathcal{G}} + \epsilon},
    \qquad i \in \mathcal{G}
\end{equation}

This restores zero mean among good samples, eliminating the
bias~\eqref{eq:bias}. Bad samples are excluded from the retain gradient
entirely (equivalent to advantage $= 0$). This matches the approach taken in
\texttt{rl-gradient-routing}.

\paragraph{Loss scaling.}
When computing the retain pass loss over only good samples, the loss should be
scaled by $(G - K) / G$ so that each good sample contributes $1/G$ to the
gradient, independent of its group's filtering rate. Without this scaling,
groups with heavy filtering (small $G - K$) would have each remaining good
sample contribute disproportionately more.

\paragraph{Small-group pathology.}
The main weakness of renormalization is statistical: when $K$ is large (many bad
samples), the good subset $\mathcal{G}$ is small. Z-scoring $n$ samples always
produces advantages with unit variance regardless of $n$, so two good samples
generate advantages of approximately $\pm 1$ --- maximal confidence from minimal
evidence. This problem worsens over training as $K$ grows
(Section~\ref{sec:retain}), degrading gradient quality precisely when the retain
adapter most needs to maintain its performance.

\subsection{Approach 2: Conditioned Reward Penalty}

Penalize bad samples proportionally to how far they exceed the good-sample mean:
\begin{equation}
    r_i^{\text{retain}} = r_i - \lambda \cdot (r_i - \bar{r}_{\mathcal{G}}),
    \qquad i \in \mathcal{B}
\end{equation}
where $\lambda$ is a dimensionless parameter. Then compute retain advantages
over the full group via standard GRPO normalization of
$\{r_i^{\text{retain}}\}_{i=1}^{G}$.

At $\lambda = 0$, this recovers the status quo.  At $\lambda = 1$, bad samples
are set to the good-sample mean, so their retain advantage is approximately zero
--- close to renormalization, but not identical.

\paragraph{Std deflation artifact.}
At $\lambda = 1$, all $K$ bad samples sit at $\bar{r}_{\mathcal{G}}$,
contributing zero variance. The group std is deflated:
\begin{equation}
    \sigma_{r^{\text{retain}}} = \sqrt{\frac{G - K - 1}{G - 1}}
        \cdot \sigma_{\mathcal{G}}
\end{equation}
For $K = 4, G = 16$: $\sigma_{r^{\text{retain}}} \approx 0.86\,
\sigma_{\mathcal{G}}$. This inflates good-sample advantages by ${\sim}17\%$
relative to pure renormalization, with the amplification increasing with $K$.
Groups with more filtering get \emph{more} influence --- the opposite of the
scaling principle in Section~\ref{sec:renorm}.

\paragraph{Reward conflation.}
The gap $r_i - \bar{r}_{\mathcal{G}}$ reflects the total reward, which includes
both the retain-task component and the RH component. A bad sample that also
performs well on the retain task receives a \emph{larger} penalty than one that
does not. This conflation is undesirable: the penalty should reflect the fact
that the sample was flagged as RH, not how much total reward it received.

\subsection{Approach 3: Manual Advantage Assignment}

Renormalize good-sample advantages via~\eqref{eq:renorm}, then directly assign
bad samples a fixed negative advantage $-\alpha$:
\begin{equation}
    \hat{A}_i^{\text{retain}} =
    \begin{cases}
        (r_i - \bar{r}_{\mathcal{G}}) / (\sigma_{\mathcal{G}} + \epsilon)
            & i \in \mathcal{G} \\
        -\alpha & i \in \mathcal{B}
    \end{cases}
\end{equation}

This cleanly decouples the good-sample signal from the avoidance signal.
However, it \textbf{subverts GRPO's normalization properties}:
\begin{itemize}[nosep]
    \item \textbf{No adaptive temperature}: a fixed $-\alpha$ does not respond
        to the reward distribution. When rewards are tightly clustered (and GRPO
        would amplify advantages via small $\sigma_r$), the avoidance signal
        remains at $\alpha$. When rewards are spread (and GRPO would compress),
        $\alpha$ is unchanged.
    \item \textbf{No per-group calibration}: ``easy'' prompts (tight rewards)
        and ``hard'' prompts (wide spread) receive the same avoidance magnitude,
        though GRPO would normally differentiate them.
    \item \textbf{No tracking over training}: as the policy evolves and the
        reward distribution shifts, $\alpha$ becomes stale.
\end{itemize}

These are exactly the properties that make GRPO stable in practice. Bypassing
them reintroduces manual tuning problems that group normalization was designed
to solve.

\subsection{Approach 4: Fixed Reward Penalty (Recommended)}\label{sec:penalty}

Apply a fixed penalty $P > 0$ (in reward-units) to bad samples' rewards, then
compute retain advantages via standard GRPO normalization over the full group:
\begin{equation}\label{eq:penalty}
    r_i^{\text{retain}} =
    \begin{cases}
        r_i & i \in \mathcal{G} \\
        r_i - P & i \in \mathcal{B}
    \end{cases}
\end{equation}
\begin{equation}\label{eq:penalty-advantage}
    \hat{A}_i^{\text{retain}}
    = \frac{r_i^{\text{retain}} - \bar{r}^{\text{retain}}}
           {\sigma_{r^{\text{retain}}} + \epsilon}
\end{equation}
where $\bar{r}^{\text{retain}}$ and $\sigma_{r^{\text{retain}}}$ are computed
over all $G$ penalized rewards. The forget adapter continues to use
unpenalized advantages from~\eqref{eq:grpo-advantage}.

\paragraph{Why this works.}
\begin{enumerate}[nosep]
    \item \textbf{Preserves GRPO normalization}: the penalty modifies the reward
        landscape; GRPO's adaptive temperature, per-group calibration, and
        scale invariance operate normally on the penalized rewards.
    \item \textbf{No small-group pathology}: all $G$ samples participate in the
        mean and std computation. No reduced group size, no noisy z-scores from
        2--3 samples.
    \item \textbf{Active avoidance}: bad samples receive negative retain
        advantages (for sufficiently large $P$), giving the retain adapter an
        explicit ``push away from RH'' signal.
    \item \textbf{Self-regulating}: if $P$ is too large, $\sigma_{r^{\text{retain}}}$
        increases (the gap between good and penalized-bad widens the distribution),
        which compresses advantages. The system absorbs the excess rather than
        blowing up.
    \item \textbf{No reward conflation}: the penalty is independent of the
        reward value. A bad sample is penalized by $P$ regardless of whether
        it also performed well on the retain task.
\end{enumerate}

\paragraph{Comparison with renormalization.}
At $P = 0$, this recovers the status quo (negative mean bias). As $P$ increases,
bad samples' advantages decrease continuously. At some $P^*$ (dependent on the
group's reward distribution), bad samples' mean retain advantage reaches zero,
approximating renormalization. Beyond $P^*$, bad samples have negative retain
advantage (active avoidance). Unlike the conditioned penalty
($\lambda$-parameterization), this does not exactly replicate renormalization at
any finite $P$, but the qualitative behavior is a smooth continuum from
``no correction'' through ``filtering'' to ``active avoidance.''

\paragraph{The penalty parameter $P$.}
$P$ is in reward-units. This is natural: it lives alongside other reward-space
hyperparameters (\texttt{scale}, \texttt{max\_reward}, etc.) that the
experimenter already configures in YAML. The GRPO normalization makes the system
robust to the exact choice of $P$ --- the penalty shifts the reward landscape,
and the normalization adapts the gradient scale accordingly.

\paragraph{Dynamics throughout training.}
\begin{itemize}[nosep]
    \item \textbf{Early} (few bad samples, $K \approx 0$): the penalty affects
        few samples. Retain advantages $\approx$ standard GRPO. Both adapters
        learn normally.
    \item \textbf{Mid} ($K \approx G/4$): penalized bad samples have clearly
        negative retain advantage. The retain adapter receives both ``improve
        good'' and ``avoid bad'' signals from the full group of $G$ samples.
    \item \textbf{Late} ($K \approx G/2$ or more): many penalized samples shift
        $\bar{r}^{\text{retain}}$ down significantly, giving all good samples
        clearly positive retain advantage. Strong ``pursue good'' signal despite
        good samples being a minority. No small-group noise because all $G$
        samples contribute to normalization.
\end{itemize}

Compare this to renormalization in the late regime: good samples are a small
minority, z-scored among themselves with high variance and low statistical power.
The fixed penalty avoids this entirely.

%% ====================================================================
\section{Summary}
%% ====================================================================

\begin{table}[h]
\centering
\begin{tabular}{@{}lcccc@{}}
\toprule
 & Retain bias & Forget bias & Renormalize? & Modify? \\
\midrule
Retain adapter & $\bar{A}_{\mathcal{G}} < 0$ (harmful)
    & --- & ---
    & \textbf{Yes} (penalty) \\
Forget adapter & --- & $\bar{A}_{\mathcal{B}} > 0$ (beneficial)
    & \textbf{No}
    & \textbf{No} \\
\bottomrule
\end{tabular}
\caption{Asymmetric treatment: correct the retain adapter's bias via a fixed
    reward penalty; leave the forget adapter's full-group advantages unchanged.}
\end{table}

\begin{table}[h]
\centering
\small
\begin{tabular}{@{}lcccc@{}}
\toprule
Property & Renormalization & Conditioned & Manual & \textbf{Fixed penalty} \\
 & (\S\ref{sec:renorm}) & penalty & advantage & (\S\ref{sec:penalty}) \\
\midrule
Fixes negative mean bias    & Yes & Yes & Yes & \textbf{Yes} \\
GRPO normalization preserved & Partially\textsuperscript{a}
    & Partially\textsuperscript{b} & No & \textbf{Yes} \\
No small-group pathology     & No  & Yes & No  & \textbf{Yes} \\
Active avoidance signal      & No  & Yes ($\lambda > 1$) & Yes & \textbf{Yes} \\
No reward conflation         & N/A & No  & N/A & \textbf{Yes} \\
New hyperparameters          & 0   & 1 ($\lambda$) & 1 ($\alpha$) & \textbf{1} ($P$) \\
\bottomrule
\end{tabular}
\caption{Comparison of corrective approaches.
    \textsuperscript{a}Std computed over reduced group.
    \textsuperscript{b}Std deflated by dead-at-mean samples.}
\end{table}

\paragraph{Recommended design.}
Apply a fixed reward penalty $P$ to RH-flagged samples' retain rewards
(Eq.~\ref{eq:penalty}), then compute retain advantages via standard GRPO
normalization over the full group (Eq.~\ref{eq:penalty-advantage}). Leave
forget-adapter advantages unchanged. This preserves GRPO's adaptive properties,
avoids small-group pathologies, and gives the retain adapter a principled
avoidance signal with a single interpretable hyperparameter.

\subsection{Implementation Site}

The intervention point is the \texttt{\_generate\_and\_score\_completions}
override in \texttt{SampleGRPOTrainer}. After the base class computes
advantages and the detector labels samples:
\begin{enumerate}[nosep]
    \item Reconstruct raw rewards from \texttt{CachedReward.\_last\_scores}
        on each component of the \texttt{CombinedReward}.
    \item Apply penalty: $r_i^{\text{retain}} = r_i - P \cdot
        \mathbf{1}[i \in \mathcal{B}]$.
    \item Reshape to $(-1, G)$ (group structure is intact pre-shuffle),
        compute per-group mean and std, normalize to get
        \texttt{retain\_advantages}.
    \item Store in the output dict. The tensor survives TRL's shuffle/split
        pipeline (applied uniformly to all dict entries).
\end{enumerate}

In \texttt{training\_step}, swap \texttt{inputs["advantages"]} with
\texttt{inputs["retain\_advantages"]} before the retain pass.
\texttt{compute\_loss} reads \texttt{inputs["advantages"]} passively ---
no changes required.

\end{document}
